% Q2V4: Parallel RLC Natural Response, Underdamped — Exam diagram
% Pre-charged capacitor V0, switch closes at t=0 connecting R||L
\documentclass[border=10pt]{standalone}
\usepackage[american voltages, american currents]{circuitikz}
\usepackage{amsmath}
\renewcommand{\familydefault}{\sfdefault}

\begin{document}
\begin{circuitikz}[line width=0.8pt, every node/.style={font=\sffamily}]

% === Pre-charged capacitor (left branch) ===
\draw (0,0) to[C, l=$C$, v=$V_0$] (0,4);

% === Closing switch from capacitor top to parallel branches ===
\draw (0,4) to[closing switch] (3,4);

% Switch labels (separate nodes, spaced vertically)
\node[above, font=\sffamily\itshape\small] at (1.5,4.6) {SW1};
\node[above, font=\sffamily\footnotesize] at (1.5,4.2) {$t{=}0$ (closes)};

% === Top wire ===
\draw (3,4) -- (6,4);

% === Bottom wire ===
\draw (0,0) -- (6,0);

% === Parallel R and L branches ===
\draw (3,4) to[R, l=$R$, i>^=$i_R$] (3,0);
\draw (6,4) to[L, l=$L$, i>^=$i_L$] (6,0);

% === Voltage across combination ===
\draw (7,4) to[open, v^=$v(t)$] (7,0);
\draw (6,4) -- (7,4);
\draw (6,0) -- (7,0);

% === Junction dots ===
\fill (0,0) circle (2pt); \fill (0,4) circle (2pt);
\fill (3,0) circle (2pt); \fill (3,4) circle (2pt);
\fill (6,0) circle (2pt); \fill (6,4) circle (2pt);

% === Initial condition label (below capacitor, clear of C label) ===
\node[left, font=\sffamily\small, text=blue!60!black] at (-0.3,0.8) {$v(0^-)=V_0$};

\end{circuitikz}
\end{document}
